%===================================================================%
%===================================================================%
\chapter{Introduction}\label{ch:introduction}
\markboth{Introduction}{}

%===================================================================%
%===================================================================%
%--- Approx Length ---%
% ~5 pages
%---------------------%

%===================================================================%
\section{Motivation and Context} % Hide Caption later !!!
%===================================================================%

%===================================================================%
\section{Problem Statement and Objectives} % Hide Caption later !!!
%===================================================================%
%---- Problem Statement ----%
Autonomous sorting of textiles from mixed waste streams is a challenging manipulation problem: garments such as T-shirts are deformable, exhibit highly variable shapes, and interact with their environment through complex contact dynamics. 
In industrial recycling contexts, reliable textile extraction is a key prerequisite for downstream steps such as material recovery or further sorting by condition (e.g., damaged vs.\ reusable). 
Classical, model-based manipulation pipelines typically rely on accurate state estimation and carefully tuned controllers, which becomes difficult when both the manipulated object and the contact conditions change continuously.

This thesis addresses the problem of learning a robust control policy for a custom tendon-driven robot to selectively retrieve a T-shirt from a conveyor transporting heterogeneous waste items and to deposit it into a designated bin. 
The employed robotic system is realized in NVIDIA Isaac~Sim and consists of three main subsystems: (i) a central tendon-driven arm section derived from an existing tendon-robot design, (ii) an articulated base that provides additional linear actuation in the $y$- and $z$-directions to expand the reachable workspace and to improve access to items on the conveyor, and (iii) a Robotiq gripper as the end-effector. 
The control policy is trained in IsaacLab using \ac{PPO} implemented with \texttt{skrl}. 
The overall learning objective is defined as a selective pick-and-throw behavior: \emph{grasp the T-shirt, lift it from the moving belt, and reliably eject it into the target bin while leaving non-target waste items behind}.

The scope of this work focuses on the learning-based manipulation stack in simulation and its preparation for sim-to-real transfer. 
Perception is treated as an enabling component (e.g., via ground-truth labels in simulation or externally provided detections) rather than the primary research contribution. 
In the broader project context, this selective retrieval capability forms the foundation for future extensions towards \emph{condition-based} T-shirt sorting, where policies must not only identify garments but also manipulate and present them in ways that enable inspection and quality assessment.

%---- Objectives ----%
The thesis pursues the following objectives:

\begin{enumerate}
  \item \textbf{High-fidelity simulation environment:}
  Develop a reproducible Isaac~Sim scene that models a conveyor-based waste stream, a target bin, and representative waste items including a deformable T-shirt, with randomized initial conditions and physical parameters to support robust learning.

  \item \textbf{Robot model integration and validation:}
  Integrate the custom robot (tendon-driven arm segment, linear $y/z$ base actuation, and Robotiq gripper) into IsaacLab, and verify the correctness of kinematics, actuation interfaces, joint limits, and controller settings required for stable simulation-based training.

  \item \textbf{\acs{MDP} formulation of the manipulation task:}
  Formulate the problem as an \ac{MDP} within the IsaacLab workflow by specifying observation space, action space, reward design, and termination criteria. The formulation must explicitly encode \emph{selectivity} (T-shirt vs.\ non-target items) and \emph{successful disposal} into the bin.

  \item \textbf{Curriculum from basic skills to the final task:}
  Implement a staged training procedure that progresses from foundational behaviors (e.g., reaching, grasping, simple pick-and-place) to the full conveyor-based selective pick-and-throw task, enabling sample-efficient learning and reducing reward sparsity.

  \item \textbf{Policy training and quantitative evaluation:}
  Train \ac{PPO} policies using \texttt{skrl} and evaluate performance with metrics such as task success rate, false-pick rate (non-target retrieval), missed-pick rate (target left on conveyor), bin placement success, cycle time, and robustness under domain randomization (e.g., belt speed, friction, mass, cloth stiffness).

  \item \textbf{Sim-to-real readiness and integration pathway:}
  Assess transfer robustness via noise and parameter randomization and define an executable integration pathway towards real-system deployment (e.g., \ac{ROS}~2 interfaces, logging, and safety constraints), identifying remaining gaps that must be closed for real-world operation.

  \item \textbf{Foundation for condition-based sorting:}
  Derive requirements for extending the learned manipulation policy towards condition-based T-shirt sorting, including which additional observations, interaction primitives (e.g., stretching/presentation), and evaluation criteria would be needed in subsequent work.
\end{enumerate}

%===================================================================%
\section{Thesis Outline} % Hide Caption later !!!
%===================================================================%

This thesis is structured as follows.
Chapter~\ref{ch:} introduces the fundamentals required for this work, covering reinforcement learning for robotic manipulation with an emphasis on \ac{PPO}, deformable-object (cloth) manipulation challenges, and the simulation tools used in this project (Isaac~Sim and IsaacLab).
Chapter~\ref{ch:system} describes the simulated system and task setup in detail: the conveyor scenario with mixed waste items, the selective retrieval objective, and the custom robot model consisting of the tendon-driven arm segment, the additional linear $y/z$ base actuation, and the Robotiq gripper.

Chapter~\ref{ch:method} formulates the learning problem as an \ac{MDP} and explains the design of observations, actions, reward terms, and termination conditions to realize selective pick-and-throw behavior (retrieve only the T-shirt and deposit it into a bin).
Chapter~\ref{ch:training} presents the training procedure in IsaacLab using \texttt{skrl} \ac{PPO}, including curriculum design from simpler sub-skills to the full conveyor task and the use of randomization to improve robustness.

Chapter~\ref{ch:experiments} defines the experimental protocol and evaluation metrics, focusing on sorting-relevant criteria such as success rate (T-shirt in bin), false retrievals of non-target objects, misses of the target, and robustness under variations of object and conveyor properties.
Chapter~\ref{ch:results} reports the quantitative and qualitative results of the learned policies and analyzes typical success behaviors and failure cases.

Finally, Chapter~\ref{ch:discussion} discusses the implications and limitations of the approach, including simulation fidelity and sim-to-real considerations, and Chapter~\ref{ch:conclusion} concludes the thesis and outlines future work towards condition-based T-shirt sorting, where retrieval is extended to classification and handling based on garment condition.